\runninghead{The Sunday and Its Texts}
\section{The Resumed Fifth Sunday after the Epiphany and Its Texts}
\label{sec:opening}

The Gospel of this Mass is a parable that Jesus tells the crowds by the sea.
A man sows good seed in his field. ``While men were asleep, his enemy came and
oversowed cockle among the wheat and went his way.'' Nothing shows until the
blade has sprung up and borne fruit; then the cockle appears too. The servants
of the householder come to him with two questions: ``Sir, didst thou not sow
good seed in thy field? Whence then hath it cockle?'' He answers, ``An enemy
hath done this.'' They ask whether they should go and gather it up, and he
refuses: ``No, lest perhaps gathering up the cockle, you root up the wheat also
together with it. Suffer both to grow until the harvest, and in the time of
the harvest I will say to the reapers: Gather up first the cockle, and bind it
into bundles to burn, but the wheat gather ye into my barn'' (Mt 13:24--30).
The Gospel stops there. The explanation that Jesus gives to the disciples
afterwards in the house, that the field is the world, the enemy the devil, the
harvest the end of the world and the reapers the angels (Mt 13:36--43), is
not read.

The Epistle, from the third chapter of Colossians, gives a list of things to
put on: ``the bowels of mercy, benignity, humility, modesty, patience: bearing
with one another and forgiving one another, if any have a complaint against
another. Even as the Lord hath forgiven you, so do you also.'' Above all these
is charity, ``the bond of perfection.'' St Paul goes on: ``let the peace of
Christ rejoice in your hearts \dots\ Let the word of Christ dwell in you
abundantly,'' teaching and admonishing ``in psalms, hymns and spiritual
canticles, singing in grace in your hearts to God,'' and doing all things in
the name of the Lord Jesus (Col 3:12--17).

Two of the five chants are promises, God's through Jeremias and Christ's in
Mark; the other three are the people's own words, thanks in the Gradual and a
cry from the depths in the Alleluia and the Offertory. The Introit gives God's
words through Jeremias to the exiles in Babylon: ``I think
thoughts of peace, and not of affliction: you shall call upon Me, and I will
hear you; and I will bring back your captivity from all places,'' with the
psalm verse ``Lord, Thou hast blessed Thy land: Thou hast turned away the
captivity of Jacob'' (Jer 29:11--12, 14; Ps 84:2). The Gradual thanks God
for a deliverance: ``Thou hast delivered us, O Lord, from them that afflict
us: and hast put them to shame that hate us'' (Ps 43:8--9). The Alleluia and
then the Offertory sing the opening of Psalm~129, ``From the depths I have
cried to Thee, O Lord: Lord, hear my prayer,'' and the Offertory repeats its
first line. The Communion is the Lord's promise in Mark: ``whatsoever you ask
when you pray, believe that you shall receive, and it shall be done unto
you'' (Mk 11:24). The English of the chants is that of the Lasance hand missal;
the lessons are given in the Douay--Rheims.

The three prayers of the priest ask for protection, pardon and fulfilment. The
Collect asks God to keep his family with continual mercy, because it leans on
the hope of heavenly grace alone. The Secret offers \latin{hostias
placationis}, offerings of propitiation, and asks that God absolve our sins
and direct our wavering hearts. The Postcommunion asks that we may receive
the effect of the salvation whose pledge we have received in these mysteries.

The Missal prints this Mass in the Proper of Time after the Twenty-third
Sunday after Pentecost, under the heading \latin{Dominica quinta quae
superfuit post Epiphaniam}, the fifth Sunday ``left over'' after the Epiphany,
of the second class (printed pp.~420--421, nos.~1749--1758). Its Collect,
Epistle, Gospel, Secret and Postcommunion are those of the Fifth Sunday after
the Epiphany; its Introit, Gradual, Alleluia, Offertory and Communion are those
of the Twenty-third Sunday after Pentecost, as RGMR~298 directs for every
Sunday after the Epiphany that is resumed between the twenty-third and the
last Sunday after Pentecost. In 2026 Septuagesima fell on 1~February and left
no room for the Fourth, Fifth and Sixth Sundays after the Epiphany. There are
twenty-six Sundays after Pentecost this year, and in that case RG~18~b sets
the Fifth Sunday after the Epiphany in the twenty-fourth place, which falls on
8~November, and the Sixth on the following Sunday; the Fourth has no place.
The calendar commemoration of the Four Crowned Martyrs on 8~November is not
made, because a Sunday of the second class admits only the commemoration of a
feast of the second class (RG~111~b). The Mass is said in green, with the
Gloria and the Credo, and with the Preface of the Most Holy Trinity, which the
Missal directs after the Secret (RGMR~494~b).

The texts raise three questions, and each of the three interpretations that
follow takes up one. The first is why the householder forbids his servants to
pull up the cockle, and what that refusal asks of those who hear it. St
Augustine, preaching on this Gospel, gives the servants the reason ``ye are
such as may be deceived'' and tells the cockle that ``they who to-day are tares,
may to-morrow be wheat''; St Jerome finds the reason in the plant, which in the
blade cannot be told from wheat; St John Chrysostom says the command forbids
the killing of heretics; and St Rabanus Maurus reads the householder's
patience as a pattern for bearing injuries. In the Epistle Paul asks the
Colossians to bear with one another and to forgive as the Lord forgave them,
and the first interpretation finds in those requests the servants' work while
both grow.

The second question is what the chants make of the people who sing them, and
why a God who has already promised to bring his people back still tells them
to call upon him. St Jerome, on the verses of Jeremias the Introit takes, says
that God could have done what he promised without the captives' prayer, and
urges them to pray so that they may receive it. On the psalm sung at the
Alleluia and the Offertory, St Hilary of Poitiers says that it is God's nature
to hear and the part of faith to pray to be heard. St Augustine reads the
captivity of the Introit's psalm as the captivity of sin, the Gradual as the
voice of the martyrs, and the depths of Psalm~129 as this mortal life.

The third question starts from four words of the Gospel, ``while men were
asleep.'' The parable does not blame the sleepers, and the Lord's explanation
does not mention them. St John Chrysostom says the danger hangs over the
rulers of the Church ``and not the rulers only, but the subjects too,'' and
that the sleep of the moral faculty can be avoided. He and St Thomas Aquinas,
commenting on the Epistle, say that the word of Christ must dwell in a man
richly and that he must sing to God in his heart and not only with his lips.
In the third interpretation that indwelling word and that singing are the
watch Chrysostom asks for.

\begin{maptable}
Introit &
\latin{Dicit Dominus: Ego cogito cogitationes pacis}; Jer 29:11--12, 14;
verse Ps 84:2 (85:2) &
The Lord thinks thoughts of peace and not of affliction; you shall call upon
him and he will hear; he will bring back your captivity from all places; the
psalm thanks him for turning away the captivity of Jacob.\\
Collect &
\latin{Familiam tuam, quaesumus, Domine, continua pietate custodi} &
That God keep his family with continual mercy, so that, leaning on the hope of
heavenly grace alone, it may always be defended by his protection.\\
Epistle &
Col 3:12--17 &
Put on mercy, kindness, humility, modesty and patience; bear with and forgive
one another as the Lord forgave you; above all have charity; let Christ's
peace rejoice in your hearts and his word dwell in you abundantly; sing to God
in your hearts; do all in the name of the Lord Jesus.\\
Gradual &
\latin{Liberasti nos, Domine, ex affligentibus nos}; Ps 43:8--9 (44:8--9) &
God has delivered us from those who afflict us and put to shame those who
hate us; in God we will glory all the day and praise his name for ever.\\
Alleluia &
\latin{De profundis clamavi ad te, Domine}; Ps 129:1--2 (130:1--2) &
From the depths I have cried to the Lord; Lord, hear my prayer.\\
Gospel &
Mt 13:24--30 &
An enemy sows cockle among the wheat while men sleep; the householder forbids
his servants to gather it, lest they root up the wheat with it, and leaves
both to grow until the harvest.\\
Offertory &
\latin{De profundis clamavi ad te, Domine}; Ps 129:1--2 (130:1--2) &
The cry from the depths sung again, its first line repeated.\\
Secret &
\latin{Hostias tibi, Domine, placationis offerimus} &
Offerings of propitiation; that God mercifully absolve our sins and direct
our wavering hearts.\\
Communion &
\latin{Amen dico vobis, quidquid orantes petitis}; Mk 11:24 &
Whatever you ask when you pray, believe that you shall receive, and it shall
be done for you.\\
Postcommunion &
\latin{Quaesumus, omnipotens Deus: ut illius salutaris capiamus effectum} &
That we may receive the effect of the salvation whose pledge we have received
in these mysteries.\\
\end{maptable}
